\subsubsection{任务启用与资源先后的显式连接}
式\eqref{eq:q2_disjoint}表达同一资源上的析取关系。为展示它如何转化为优化约束，对不同任务$r,r'$和同型飞机$u$，引入$z^U_{rr'u}\in\{0,1\}$，表示两项任务同时分给$u$时是否由$r$先执行。在有限候选模型中，取覆盖所考虑时间域及最长占用长度的常数$M_U$，可写为
\begin{equation}
\begin{aligned}
s_{r'}&\ge s_r+p_r-M_U(1-z^U_{rr'u})-M_U(2-a_{ru}-a_{r'u}),\\
s_r&\ge s_{r'}+p_{r'}-M_Uz^U_{rr'u}-M_U(2-a_{ru}-a_{r'u}).
\end{aligned}
\label{eq:q2_plane_disjunction}
\end{equation}
当两项任务都使用$u$时，两式只保留选定方向的先后约束；至少一项没有分给$u$时，两个不等式均被解除。候选时间域可由已知可行日程给定，$M_U$应由实际上界推得，而不是使用没有量纲依据的极大数。电池关系采用独立先后变量，并将飞机占用长度$p_r$替换为$p_r+c_r$，同时为充电结束时刻保留相应上界。

这一显式表达说明：飞机与电池可以分别有不同的前后任务，同一开始时刻$s_r$则必须满足两条资源链的释放条件。它与事件排程对应同一个物理模型。前者方便解释完整约束，后者在大批候选之间构造可行日程更直接；使用事件排程并不意味着忽略了两类资源冲突。

\subsubsection{分层目标与搜索评分的分工}
硬截止先作为可行性约束，随后按$(L_w,T_{\max},E,N)$比较候选。该顺序可由逐层固定最优值实现，也可直接用于可行方案档案的字典序比较。其含义是，前一层取得改善时，不用后一层的变化抵消它；所有指标仍分别保留原始单位和数值。

搜索阶段采用的标量评分用于引导探索，允许暂时接受较差当前状态以离开局部结构。正式最好解从经过完整检查的可行档案中选择，始终采用既定目标顺序。因此，搜索评分中的能耗系数并不是将1 kWh与若干秒等价为社会成本。论文通过具体方案和独立指标呈现取舍，使算法参数与救援评价保持分工。

少飞一个架次也未必缩短工期。拆分任务可能降低单次载荷、提高并行性，却增加准备、交接和资源接续；合并任务则可能节省固定成本，却把关键箱子的交付推迟。候选变化必须经过完整任务计算和双资源时间表生成，才能评价系统效应。
